% !PRINTED 641
\weberpage{641}
\begin{center}
\S.\,176.\\[4pt]
\gesperrt{The real field $H_m$ contained in $\Omega_m$.}
\end{center}

Every number of the field $\Omega_m$ passes by the substitution $(r, r^{-1})$ into the conjugate imaginary number, which is also obtained by the interchange $(i, -i)$. The product of two such conjugate imaginary numbers is the \gesperrt{square of the absolute value} of each one of them.

A number in $\Omega_m$ is real when it remains unchanged under the interchange $(r, r^{-1})$, and only under this assumption. Every such number can be expressed rationally in terms of the \gesperrt{two-term period} $r + r^{-1}$, and therefore belongs to a real field of degree $\frac12\varphi(m)$, which is a divisor of $\Omega_m$. This field we will call the \gesperrt{real field contained in $\Omega_m$} and denote by $H_m$ (although still other real fields, namely all divisors of $H_m$, are contained in $\Omega_m$ as well).

The $\frac12\varphi(m) = \frac12\mu$ numbers
\weq{1}{1,\ r+r^{-1},\ r^2+r^{-2},\ \dots,\ r^{\frac12\mu-1}+r^{-\frac12\mu+1}}
form a basis of the integers of the field $H_m$. For according to \S.\,170, (7),
\weq{2}{\omega = x_0+x_1r+x_2r^2+\cdots+x_{\frac12\mu-1}r^{\frac12\mu-1}+x_{\frac12\mu}r^{\frac12\mu} +x_{-1}r^{-1}+x_{-2}r^{-2}+\cdots+x_{-\frac12\mu+1}r^{-\frac12\mu+1}}
is then and only then an integer of the field $\Omega_m$, when the coefficients $x_0, x_1, \dots$ are integral rational numbers.

The condition for this number to belong to the field $H_m$ is obtained by carrying out the substitution $(r, r^{-1})$ and setting the difference $=0$. If we further divide by $r-r^{-1}$, there results in this way:
\[
(x_1-x_{-1})+(x_2-x_{-2})(r+r^{-1})+\cdots
+(x_{\frac12\mu-1}-x_{-\frac12\mu+1})\frac{r^{\frac12\mu-1}-r^{-\frac12\mu+1}}{r-r^{-1}}
+x_{\frac12\mu}\frac{r^{\frac12\mu}-r^{-\frac12\mu}}{r-r^{-1}}=0.
\]
The divisions can here be carried out, and if we then multiply by $r^{\frac12\mu-1}$, there results for $r$ a rational equation of degree lower than $\mu$, which can only
%CONT>>
